\documentclass[10pt,a4paper]{amsart}
\usepackage{amsmath,amssymb,amsthm,graphicx}
% Retained minimal preamble for the Stacks Project chapter more-algebra.tex (More on Algebra).
% The chapter file pinned-original-source.tex carries only \input{preamble} on its first line, so no
% preamble text lives inside the chapter itself.  The Stacks Project uses one single shared preamble for
% every chapter; the snapshot stages it identically as sources/stacks-chow/preamble.tex and
% sources/stacks-etale-cohomology/preamble.tex.  The block below is copied mechanically, line for line,
% from sources/stacks-chow/preamble.tex lines 14-18 (the Xy-pic setup) and 154-217 (the native theorem
% declarations together with the mathematical macros).  Nothing was added, removed or redefined; the
% wrapper class and the amsmath/amssymb/amsthm/graphicx base packages are supplied by the assembler.
% Chapter: more-algebra.tex -- More on Algebra.
% For commutative diagrams we use Xy-pic
\usepackage[all]{xy}

% We use 2cell for 2-commutative diagrams.
\xyoption{2cell}

% Theorem environments.
%
\theoremstyle{plain}
\newtheorem{theorem}[subsection]{Theorem}
\newtheorem{proposition}[subsection]{Proposition}
\newtheorem{lemma}[subsection]{Lemma}

\theoremstyle{definition}
\newtheorem{definition}[subsection]{Definition}
\newtheorem{example}[subsection]{Example}
\newtheorem{exercise}[subsection]{Exercise}
\newtheorem{situation}[subsection]{Situation}

\theoremstyle{remark}
\newtheorem{remark}[subsection]{Remark}
\newtheorem{remarks}[subsection]{Remarks}

\numberwithin{equation}{subsection}

% Macros
%
\def\lim{\mathop{\mathrm{lim}}\nolimits}
\def\colim{\mathop{\mathrm{colim}}\nolimits}
\def\Spec{\mathop{\mathrm{Spec}}}
\def\Hom{\mathop{\mathrm{Hom}}\nolimits}
\def\Ext{\mathop{\mathrm{Ext}}\nolimits}
\def\SheafHom{\mathop{\mathcal{H}\!\mathit{om}}\nolimits}
\def\SheafExt{\mathop{\mathcal{E}\!\mathit{xt}}\nolimits}
\def\Sch{\mathit{Sch}}
\def\Mor{\mathop{\mathrm{Mor}}\nolimits}
\def\Ob{\mathop{\mathrm{Ob}}\nolimits}
\def\Sh{\mathop{\mathit{Sh}}\nolimits}
\def\NL{\mathop{N\!L}\nolimits}
\def\CH{\mathop{\mathrm{CH}}\nolimits}
\def\proetale{{pro\text{-}\acute{e}tale}}
\def\etale{{\acute{e}tale}}
\def\QCoh{\mathit{QCoh}}
\def\Ker{\mathop{\mathrm{Ker}}}
\def\Im{\mathop{\mathrm{Im}}}
\def\Coker{\mathop{\mathrm{Coker}}}
\def\Coim{\mathop{\mathrm{Coim}}}

% Boxtimes
%
\DeclareMathSymbol{\boxtimes}{\mathbin}{AMSa}{"02}

%
% Macros for moduli stacks/spaces
%
\def\QCohstack{\mathcal{QC}\!\mathit{oh}}
\def\Cohstack{\mathcal{C}\!\mathit{oh}}
\def\Spacesstack{\mathcal{S}\!\mathit{paces}}
\def\Quotfunctor{\mathrm{Quot}}
\def\Hilbfunctor{\mathrm{Hilb}}
\def\Curvesstack{\mathcal{C}\!\mathit{urves}}
\def\Polarizedstack{\mathcal{P}\!\mathit{olarized}}
\def\Complexesstack{\mathcal{C}\!\mathit{omplexes}}
% \Pic is the operator that assigns to X its picard group, usage \Pic(X)
% \Picardstack_{X/B} denotes the Picard stack of X over B
% \Picardfunctor_{X/B} denotes the Picard functor of X over B
\def\Pic{\mathop{\mathrm{Pic}}\nolimits}
\def\Picardstack{\mathcal{P}\!\mathit{ic}}
\def\Picardfunctor{\mathrm{Pic}}
\def\Deformationcategory{\mathcal{D}\!\mathit{ef}}

% --- appended 2026-10-12 (Chief): the project own definitions of the reference/slogan/history
% display environments, copied mechanically from the shared Stacks preamble lines 8-12
% (\usepackage{verbatim} plus the three \newenvironment lines, which the chapter preamble keeps
% verbatim from the project preamble). Nothing added, removed or redefined: these environments
% carry no mathematical content and the project itself typesets them as comments. Without these
% three lines a counted statement whose own text uses them fails with Environment slogan undefined.
\usepackage{verbatim}
\newenvironment{reference}{\comment}{\endcomment}
\newenvironment{slogan}{\comment}{\endcomment}
\newenvironment{history}{\comment}{\endcomment}
\usepackage{cleveref}
\crefname{originallemma}{original lemma}{original lemmas}
\Crefname{originallemma}{Original lemma}{Original lemmas}
\makeatletter
\newlabel{more-algebra-lemma-dominate-two-surjections}{{09Q8}{source}{}{sourcepointer-1}{}}
\newlabel{more-algebra-lemma-dominate-two-surjections@cref}{{[originallemma][0][]09Q8}{[1][0][]source}{}{}{}}
\makeatother
\crefname{originallemma}{original lemma}{original lemmas}
\Crefname{originallemma}{Original lemma}{Original lemmas}
\makeatletter
\newlabel{algebra-definition-cohen-ring}{{0327}{source}{}{sourcepointer-2}{}}
\newlabel{algebra-definition-cohen-ring@cref}{{[originallemma][0][]0327}{[1][0][]source}{}{}{}}
\makeatother
\crefname{originallemma}{original lemma}{original lemmas}
\Crefname{originallemma}{Original lemma}{Original lemmas}
\makeatletter
\newlabel{algebra-lemma-cohen-rings-exist}{{0328}{source}{}{sourcepointer-3}{}}
\newlabel{algebra-lemma-cohen-rings-exist@cref}{{[originallemma][0][]0328}{[1][0][]source}{}{}{}}
\makeatother
\crefname{originallemma}{original lemma}{original lemmas}
\Crefname{originallemma}{Original lemma}{Original lemmas}
\makeatletter
\newlabel{algebra-theorem-cohen-structure-theorem}{{032A}{source}{}{sourcepointer-4}{}}
\newlabel{algebra-theorem-cohen-structure-theorem@cref}{{[originallemma][0][]032A}{[1][0][]source}{}{}{}}
\makeatother
\crefname{originallemma}{original lemma}{original lemmas}
\Crefname{originallemma}{Original lemma}{Original lemmas}
\makeatletter
\newlabel{lemma-power-series-ring-over-Cohen-fs}{{07NL}{source}{}{sourcepointer-5}{}}
\newlabel{lemma-power-series-ring-over-Cohen-fs@cref}{{[originallemma][0][]07NL}{[1][0][]source}{}{}{}}
\makeatother
\crefname{originallemma}{original lemma}{original lemmas}
\Crefname{originallemma}{Original lemma}{Original lemmas}
\makeatletter
\newlabel{lemma-lift-continuous}{{07NJ}{source}{}{sourcepointer-6}{}}
\newlabel{lemma-lift-continuous@cref}{{[originallemma][0][]07NJ}{[1][0][]source}{}{}{}}
\makeatother
\title{The fibre product of two complete Noetherian local rings over a common quotient is complete Noetherian local}
\author{The Stacks Project Authors}
\date{stacks-project:more-algebra.tex:local-snapshot-20261012; GNU FDL 1.2 or later}
\begin{document}
\maketitle
\setcounter{subsection}{1}
\noindent\textit{Editorial note.} All delivered mathematical text is the author's own, copied from the pinned Stacks Project chapter file more-algebra.tex; no statement or proof step was rewritten, reordered or removed.\par\medskip
\noindent\textit{Editorial note.} The chapter file pinned-original-source.tex carries only \textbackslash{}input\{preamble\} on its first line, so no preamble text lives inside the chapter. The retained file sources/stacks-more-algebra/retained-preamble.tex is a mechanical, line-for-line copy of the project's single shared preamble as staged in the snapshot (Xy-pic setup plus theorem declarations and mathematical macros). Nothing was added, removed or redefined, and the excerpt serves only as preamble, never as a span coordinate.\par\medskip
\noindent\textit{Editorial note.} Both counted spans start at the \textbackslash{}begin line of their environment, so the statement and the author's proof are delivered whole; surrounding prose of the section is deliberately outside the counted spans and is carried as standard prerequisites instead.\par\medskip
\noindent\textit{Editorial note.} The citation fields (source\_id, source\_version, authors, source\_url, license.name, source\_title) are identical to those of every other card this lane has produced from this chapter, so the batch selector treats them as one source.\par\medskip
\noindent\textit{Allowed standard prerequisites (editorial).}\begin{itemize}
\item Existence of Cohen rings with prescribed residue field and the Cohen structure theorem. Algebra, Definition \textbackslash{}ref\{algebra-definition-cohen-ring\}, Algebra, Lemma \textbackslash{}ref\{algebra-lemma-cohen-rings-exist\}, and Algebra, Theorem \textbackslash{}ref\{algebra-theorem-cohen-structure-theorem\}, cited at proof lines 9828-9837. Provide the power series presentation of the common quotient and the characteristic split.
\item Power series rings over a Cohen ring are formally smooth, and continuous maps lift along surjections of complete local rings. Lemma \textbackslash{}ref\{lemma-power-series-ring-over-Cohen-fs\} and Lemma \textbackslash{}ref\{lemma-lift-continuous\}, cited at proof lines 9840-9842. Lift the surjections onto the common quotient to maps from the power series ring.
\end{itemize}\medskip
\ifdefined\hypertarget\hypertarget{sourcepointer-1}{}\fi
\noindent\textit{Source reference.} Original lemma 09Q8: focused-seeds/AGP/sources/stacks-snapshot-20261011/tags.txt, source lines 12276--12276. The lemma's own canonical tag, quoted so that readers can resolve the label inside the delivered text; it points at the counted statement itself and carries no proof content.\par\medskip
\ifdefined\hypertarget\hypertarget{sourcepointer-2}{}\fi
\noindent\textit{Source reference.} Original lemma 0327: focused-seeds/AGP/sources/stacks-snapshot-20261011/tags.txt, source lines 4010--4010. Navigation pointer to a statement the counted proof refers to but does not contain; nothing load-bearing is delegated to it.\par\medskip
\ifdefined\hypertarget\hypertarget{sourcepointer-3}{}\fi
\noindent\textit{Source reference.} Original lemma 0328: focused-seeds/AGP/sources/stacks-snapshot-20261011/tags.txt, source lines 4011--4011. Navigation pointer to a statement the counted proof refers to but does not contain; nothing load-bearing is delegated to it.\par\medskip
\ifdefined\hypertarget\hypertarget{sourcepointer-4}{}\fi
\noindent\textit{Source reference.} Original lemma 032A: focused-seeds/AGP/sources/stacks-snapshot-20261011/tags.txt, source lines 4013--4013. Navigation pointer to a statement the counted proof refers to but does not contain; nothing load-bearing is delegated to it.\par\medskip
\ifdefined\hypertarget\hypertarget{sourcepointer-5}{}\fi
\noindent\textit{Source reference.} Original lemma 07NL: focused-seeds/AGP/sources/stacks-snapshot-20261011/tags.txt, source lines 9766--9766. Navigation pointer to a statement the counted proof refers to but does not contain; nothing load-bearing is delegated to it.\par\medskip
\ifdefined\hypertarget\hypertarget{sourcepointer-6}{}\fi
\noindent\textit{Source reference.} Original lemma 07NJ: focused-seeds/AGP/sources/stacks-snapshot-20261011/tags.txt, source lines 9764--9764. Navigation pointer to a statement the counted proof refers to but does not contain; nothing load-bearing is delegated to it.\par\medskip
\begin{lemma}
\label{lemma-dominate-two-surjections}
Let $S \to R$ and $S' \to R$ be surjective maps of complete Noetherian local
rings. Then $S \times_R S'$ is a complete Noetherian local ring.
\end{lemma}

\begin{proof}
Let $k$ be the residue field of $R$. If the characteristic of
$k$ is $p > 0$, then we denote $\Lambda$ a Cohen ring
(Algebra, Definition \ref{algebra-definition-cohen-ring})
with residue field $k$ (Algebra, Lemma \ref{algebra-lemma-cohen-rings-exist}).
If the characteristic of $k$ is $0$ we set $\Lambda = k$.
Choose a surjection $\Lambda[[x_1, \ldots, x_n]] \to R$
(as in the Cohen structure theorem, see
Algebra, Theorem \ref{algebra-theorem-cohen-structure-theorem})
and lift this to maps $\Lambda[[x_1, \ldots, x_n]] \to S$ and
$\varphi : \Lambda[[x_1, \ldots, x_n]] \to S$ and
$\varphi' : \Lambda[[x_1, \ldots, x_n]] \to S'$ using
Lemmas \ref{lemma-power-series-ring-over-Cohen-fs} and
\ref{lemma-lift-continuous}.
Next, choose $f_1, \ldots, f_m \in S$ generating the kernel
of $S \to R$ and $f'_1, \ldots, f'_{m'} \in S'$ generating the
kernel of $S' \to R$. Then the map
$$
\Lambda[[x_1, \ldots, x_n, y_1, \ldots, y_m, z_1, \ldots, z_{m'}]]
\longrightarrow S \times_R S,
$$
which sends $x_i$ to $(\varphi(x_i), \varphi'(x_i))$ and
$y_j$ to $(f_j, 0)$ and
$z_{j'}$ to $(0, f'_j)$
is surjective. Thus $S \times_R S'$ is a quotient of a complete
local ring, whence complete.
\end{proof}

\end{document}
